\chapter{Diễn giải kết quả và khuyến nghị}

\section{Kết quả end-to-end}

Hệ thống hoàn chỉnh chuyển một cặp tweet--ảnh qua chuỗi:

\[
\text{dữ liệu}\rightarrow\text{đặc trưng}\rightarrow
\text{xác suất}\rightarrow\text{fusion}\rightarrow
\text{chính sách DSS}\rightarrow\text{hàng đợi}.
\]

Integration test sử dụng một bài đăng Hurricane Harvey về hoạt động cứu trợ.
Hệ thống dự báo đúng nhóm rescue/donation, tạo Risk Score 54,3, Priority Medium
và định tuyến Relief Team. Kiểm thử này chứng minh các module có thể phối hợp
trên một ảnh thật; nó không thay thế đánh giá thống kê trên toàn test.

Dashboard vận hành trên 2.237 dòng test để người dùng có đủ batch quan sát: 50
case High và 571 case Manual Review. Các metric khoa học chỉ tính trên 2.169
dòng canonical test; trong tập này có 49 High và 539 Manual Review. Việc tách
hai con số tránh nhầm dashboard workload với tập đánh giá đã lọc duplicate.

\section{Ý nghĩa của kết quả mô hình}

\subsection{Sàng lọc informative}

Fusion đạt F2 .9045, nhưng baseline luôn dự báo informative đã đạt .8939. Mức
tăng .0106 là nhỏ về độ lớn thực tiễn dù CI bootstrap của gain không chứa 0.
Giá trị bổ sung rõ hơn ở Accuracy (.6948 so với .6275), F1 (.8025 so với
.7711), Balanced Accuracy (.5944 so với .5000) và MCC (.3325 so với 0).

\begin{figure}[H]
\centering
\begin{tikzpicture}[font=\scriptsize]
\def\sc{7.0}
\foreach \v/\lab/\x in {0.0106/F2/0, 0.0314/F1/1.5, 0.0673/Accuracy/3.0,
  0.0944/{Bal. Acc.}/4.5, 0.3325/MCC/6.0}{
  \fill[HUSTRed!55] (\x,0) rectangle ++(0.95,\v*\sc);
  \node[above,font=\tiny] at (\x+0.47,\v*\sc) {+\v};
  \node[below,align=center,font=\tiny] at (\x+0.47,0) {\lab};
}
\draw[->] (-0.2,0) -- (7.2,0);
\node[rotate=90,anchor=south,font=\tiny] at (-0.45,1.2) {Gain so dummy};
\end{tikzpicture}
\caption{Mức tăng của Late Fusion so với baseline always-informative theo từng
metric (canonical test). F2 gần như không nhúc nhích, nhưng các metric cân bằng
(Balanced Accuracy, MCC) tăng rõ -- đây mới là giá trị thật của fusion.}
\label{fig:informative-delta}
\end{figure}

Hình~\ref{fig:informative-delta} cho thấy ngay vì sao không được khoe F2 đơn
độc: cột F2 thấp nhất, còn MCC tăng từ 0 lên .3325 mới là bằng chứng mô hình
phân biệt hai lớp hơn hẳn dự báo hằng.

Điều này có nghĩa hệ thống phân biệt hai lớp tốt hơn dự báo hằng, nhưng ưu tiên
Recall khiến nhiều not-informative vẫn bị chuyển tiếp. Trong bối cảnh cứu trợ,
đó có thể là đánh đổi chấp nhận được ở tầng sàng lọc đầu tiên, miễn là tầng sau
có hàng đợi và con người kiểm tra. Không nên mô tả informative model là bộ lọc
chính xác cao hoặc dùng nó để tự động loại bỏ mọi bài có xác suất thấp.

\subsection{Phân loại nhu cầu humanitarian}

Humanitarian fusion đạt Macro-F1 .4005, cao hơn image-only .3646, text-only
.3260 và majority baseline .0686. Mức tăng so với baseline rõ hơn task
informative vì model thực sự phân biệt nhiều lớp thay vì chỉ tận dụng tỷ lệ
lớp dương.

Tuy nhiên, kết quả không đồng đều. F1 là .6791 ở not-humanitarian, .5645 ở
other-relevant và .5363 ở rescue/donation, nhưng chỉ .0714 ở missing/found.
Do đó hệ thống có giá trị cho tổ chức hàng đợi tổng thể, chưa đủ tin cậy để một
nhãn lớp hiếm tự kích hoạt điều động sinh tử.

\begin{figure}[H]
\centering
\begin{tikzpicture}[font=\scriptsize]
\def\sc{5.5}
\foreach \v/\lab/\x in {0.0686/Majority/0, 0.3260/Text-only/1.4,
  0.3646/Image-only/2.8, 0.4005/Fusion/4.2}{
  \fill[HUSTBlue!60] (\x,0) rectangle ++(0.9,\v*\sc);
  \node[above,font=\tiny] at (\x+0.45,\v*\sc) {\v};
  \node[below,align=center] at (\x+0.45,0) {\lab};
}
\draw[->] (-0.2,0) -- (5.4,0);
\node[rotate=90,anchor=south] at (-0.4,1.1) {Macro-F1};
\end{tikzpicture}
\caption{Tiến triển Macro-F1 humanitarian qua các mốc đối chứng trên canonical
test: mỗi bước thêm thông tin đều nhích lên, fusion cao nhất.}
\label{fig:results-progression}
\end{figure}

Hình~\ref{fig:results-progression} cho thấy điều mà một bảng số khó truyền tải
ngay: khoảng cách lớn nhất là từ majority lên các mô hình thật, và fusion vượt
cả hai nhánh đơn lẻ. Người đọc không cần tự trừ các con số trong bảng.

\subsection{Giá trị của nhánh ảnh}

Image-only tốt hơn text-only ở humanitarian Macro-F1 (.3646 so với .3260), và
fusion tốt hơn cả hai. Đây là bằng chứng ảnh cung cấp thông tin bổ sung. EDA lại
cho thấy nhãn text--image bất đồng cao ở lớp khẩn cấp. Hai kết quả không mâu
thuẫn: ảnh có thể hữu ích trên trung bình nhưng vẫn không khớp văn bản ở nhiều
case cụ thể.

Kết luận vận hành là giữ ảnh như một nguồn bằng chứng độc lập, không gán cho nó
vai trò ``xác nhận chân lý''. Khi hai nhánh bất đồng mạnh, hệ thống tăng nhu cầu
xác minh thay vì tự chọn một nhánh.

\section{Ý nghĩa của Manual Review}

Với giới hạn công suất dev 25\%, ngưỡng conflict 0,54 tạo review rate 24,85\%
trên canonical test. Diễn giải bằng \emph{số ca} dễ hình dung hơn tỷ lệ: trên
2.169 dòng test, hệ thống đẩy khoảng \(539\) ca vào hàng review. Precision .7532
nghĩa là trong \(539\) ca đó có khoảng \(406\) ca thực sự bất đồng theo
annotation và khoảng \(133\) ca là báo nhầm. Recall .3295 nghĩa là hệ thống chỉ
thu hồi khoảng một phần ba \emph{toàn bộ} ca bất đồng -- phần còn lại lọt qua.
Đây là một đánh đổi công suất: muốn bắt nhiều hơn phải hạ ngưỡng và chấp nhận
hàng review dài hơn.

Đây là cơ chế \textit{triage}, không phải detector toàn diện. Nếu trung tâm có
thêm người, có thể hạ threshold để tăng Recall. Nếu hàng đợi quá tải, cần tăng
threshold hoặc ưu tiên conflict đi cùng category nghiêm trọng. Thay đổi threshold
phải được giám sát bằng cả review rate và chất lượng review, không chỉ một metric.

\section{Một số ca minh họa end-to-end}

Để thấy hệ thống ``chạy'' chứ không chỉ là bảng số, Bảng~\ref{tab:casestudies}
theo dõi ba kịch bản minh họa đi hết luồng: xác suất hai nhánh $\rightarrow$
fusion $\rightarrow$ conflict $\rightarrow$ Risk/Priority $\rightarrow$ đề xuất
$\rightarrow$ bước con người phải làm. Các số mang tính minh họa cho cơ chế, phù
hợp với các tình huống đã nêu ở Mục~\ref{sec:scenarios}.

\begin{table}[H]\centering\scriptsize
\caption{Ba ca minh họa đi hết luồng quyết định.}
\label{tab:casestudies}
\begin{tabularx}{\textwidth}{p{2.5cm}p{1cm}p{1cm}p{1cm}X}
\toprule
\textbf{Ca} & \(p_{\text{text}}\) & \(p_{\text{img}}\) & \textbf{Conflict} &
\textbf{Luồng và bước con người}\\
\midrule
A. ``Bridge collapsed, people trapped'' + ảnh hạ tầng hỏng &
0,92 & 0,88 & thấp &
Fusion cao, đồng thuận $\rightarrow$ Risk High, route Emergency Team. Người xác
minh \emph{vị trí} qua hotline trước khi điều lực lượng.\\
B. ``Urgent help needed!'' + ảnh selfie không liên quan &
0,85 & 0,30 & \textbf{cao} &
Bất đồng mạnh $\rightarrow$ giữ ưu tiên nhưng bật Manual Review. Người xem ảnh và
văn bản, đối chiếu nguồn thứ hai.\\
C. Dự báo missing/found rất tự tin nhưng lớp chỉ có ít mẫu train &
-- & -- & -- &
Xác suất cao \emph{không} đồng nghĩa đáng tin $\rightarrow$ bắt buộc xác minh
trước mọi hành động (xem lớp hiếm ở Chương~\ref{chap:mohinh}).\\
\bottomrule
\end{tabularx}
\end{table}

Điểm chung của ba ca: hệ thống chỉ \emph{đề xuất thứ tự và đội phụ trách}; mọi
hành động có hậu quả đều có một bước con người chèn vào trước.

Hình~\ref{fig:case-studies} lấy \emph{ba ca thật} từ test (ảnh và số liệu do hệ
thống sinh ra, không dàn dựng) minh họa đúng ba tình huống trên:

\begin{figure}[H]\centering
\includegraphics[width=\textwidth]{case_studies.png}
\caption{Ba ca thật từ test. A: tin đồng thuận, Risk 87,5 $\rightarrow$ High,
không cần review. B: xung đột cao (conflict 0,92) nên bật Manual Review dù vẫn
ưu tiên. C: tweet ``no missing child'' -- một câu \emph{phủ định} thật, lớp
missing/found hiếm, hệ thống chuyển review thay vì tin tuyệt đối. Ca C cho thấy
đúng điểm yếu phủ định của TF--IDF đã nêu ở Chương~\ref{chap:mohinh}.}
\label{fig:case-studies}
\end{figure}

\section{Khuyến nghị vận hành}

\begin{enumerate}[leftmargin=*]
  \item \textbf{Dùng hệ thống để sàng lọc, không tự động điều động.}
        Priority và Routing là đề xuất ban đầu; người điều phối xác minh nguồn,
        vị trí và tình trạng hiện trường trước khi phân bổ nguồn lực.
  \item \textbf{Xử lý High trước nhưng vẫn phân tầng trong High.}
        Case injured/missing, có từ khóa khẩn cấp và địa điểm rõ nên được xem
        trước case thiệt hại tài sản.
  \item \textbf{Review song song, không chặn phản ứng khẩn cấp.}
        Khi case High có conflict, Supervisor và đội phản ứng gốc cùng nhận;
        không chờ xác minh xong mới gửi tín hiệu cho Emergency Team.
  \item \textbf{Quản lý công suất Manual Review.}
        Theo dõi số case theo ca trực. Capacity thấp hơn 25\% cần tăng ngưỡng;
        capacity cao hơn có thể hạ ngưỡng sau khi kiểm tra Precision/Recall.
  \item \textbf{Không tự động hóa lớp support rất thấp.}
        Missing/found và vehicle damage cần xác minh bắt buộc cho tới khi có
        thêm dữ liệu và đánh giá ngoài mẫu.
  \item \textbf{Theo dõi drift theo sự kiện.}
        Dùng event filter, so sánh phân bố và metric. Không áp dụng nguyên
        threshold từ dữ liệu năm 2017 cho một thảm họa mới mà không hiệu chỉnh.
  \item \textbf{Phê duyệt Risk policy bằng dữ liệu nghiệp vụ.}
        Cần chuyên gia gán Priority, chi phí bỏ sót, thời gian xử lý và kết quả
        thực tế để kiểm định trọng số Risk Score.
\end{enumerate}

\section{Trả lời tám câu hỏi của bài tập}

\begin{enumerate}[leftmargin=*]
  \item \textbf{Bài toán thực tế mà nhóm giải quyết là gì?}\\
        Trong thảm họa, lượng bài đăng mạng xã hội vượt khả năng đọc thủ công,
        nội dung nhiễu và hai nguồn text--image có thể không đồng nhất. Bài toán
        là chuyển dòng dữ liệu này thành một hàng đợi có cấu trúc để giảm thời
        gian sàng lọc và hỗ trợ điều phối.

  \item \textbf{Ai là người sử dụng hệ thống?}\\
        Người dùng chính là ban điều phối và các Emergency, Relief,
        Infrastructure, Coordination Team. Supervisor xử lý case mâu thuẫn.
        Người dân không trực tiếp nhận quyết định từ model.

  \item \textbf{Hệ thống hỗ trợ những quyết định nào?}\\
        Năm quyết định gồm: bài có đáng chú ý hay không; thuộc nhóm nhu cầu nào;
        Priority Low/Medium/High; đội nào tiếp nhận; và có cần Manual Review hay
        không. Hệ thống đề xuất, con người chịu trách nhiệm quyết định cuối.

  \item \textbf{Dữ liệu nào được sử dụng?}\\
        18.082 cặp tweet--ảnh CrisisMMD v2.0 từ bảy thảm họa năm 2017, gồm nhãn
        informative, tám lớp humanitarian, nhãn riêng text/image và cờ đồng
        thuận. Nhãn informative chính thức được join theo cặp ID; duplicate
        chính xác và gần trùng được kiểm tra trước đánh giá.

  \item \textbf{Dữ liệu được xử lý và phân tích như thế nào?}\\
        Văn bản được làm sạch và biểu diễn bằng TF--IDF 2.000 chiều. Ảnh được
        chuyển thành CLIP embedding 512 chiều với trọng số CLIP giữ nguyên.
        EDA dùng thống kê, K-Means, Apriori, t-SNE và conflict analysis. Sáu
        classifier được so sánh; họ thắng được tune trên dev rồi Late Fusion
        kết hợp xác suất.

  \item \textbf{Kết quả quan trọng nhất là gì?}\\
        Informative fusion đạt Accuracy .6948, F1 .8025, F2 .9045; F2 chỉ hơn
        always-informative .0106. Humanitarian fusion đạt Macro-F1 .4005 so với
        majority .0686. Trên robust test không tune lại, F2 còn .9016 và
        Macro-F1 .3908. Lớp hiếm vẫn là điểm yếu lớn nhất.

  \item \textbf{Hành động nào được đề xuất từ kết quả?}\\
        Dùng Priority Queue, routing theo category, review song song khi
        conflict, bắt buộc xác minh lớp hiếm và hiệu chỉnh ngưỡng theo công suất.
        Giám sát kết quả theo event để phát hiện domain shift.

  \item \textbf{Hạn chế và rủi ro là gì?}\\
        Dữ liệu chỉ từ năm 2017, thiếu GPS và xác minh hiện trường, lớp hiếm rất
        nhỏ, còn tương quan trong cùng event, có duplicate, CLIP không được
        thích nghi riêng cho miền, và Priority không có ground truth. Prototype
        chưa được kiểm thử tải production hoặc tác động thực địa.
\end{enumerate}

\section{Điều đã chứng minh và chưa chứng minh}

\begin{table}[H]\centering\small
\caption{Biên giới của kết luận.}
\begin{tabularx}{\textwidth}{XX}
\toprule
\textbf{Đã có bằng chứng} & \textbf{Chưa có bằng chứng}\\
\midrule
Fusion hơn dummy về nhiều metric và hơn từng nhánh ở humanitarian &
F2 gain khoảng .01 tạo lợi ích vận hành lớn\\
Classifier đã được so sánh và tune chỉ trên dev &
Một cấu hình sẽ tổng quát hóa sang thảm họa chưa thấy\\
Conflict rule có Precision .7532 dưới capacity 25\% &
Manual Review thu hồi được mọi xung đột quan trọng\\
Routing phủ đủ tám category và qua scenario test &
Priority threshold và severity weight là tối ưu thực tế\\
Robustness/CI không đảo ngược kết luận tổng thể &
Lớp missing đủ tin cậy để tự động điều động\\
Dashboard chạy với model và cache đã khóa &
Hệ thống đáp ứng bảo mật, tải và độ trễ production\\
\bottomrule
\end{tabularx}
\end{table}

Giá trị chính của dự án nằm ở sự kết hợp giữa dự báo và quản trị quyết định:
model giảm khối lượng thông tin, còn chính sách minh bạch, capacity constraint,
scenario test và human-in-the-loop giới hạn cách kết quả được sử dụng.
